\section{Descriptive Statistics: Turning 69 Responses into a Summary}

After understanding the research variables, the next step is to use
descriptive statistics. Descriptive statistics are used to examine the
central tendency and distribution of the data before analyzing
relationships between variables.

In this study, the responses to each item were summed to produce a total
score for variables X and Y for each respondent.

Variable X consisted of 15 items, so its theoretical total score ranged
from
\[
15 \times 1 = 15
\]
to
\[
15 \times 5 = 75
\]
Therefore, the theoretical range of variable X was 15--75.

Variable Y consisted of 14 items, so its theoretical score ranged from
\[
14 \times 1 = 14
\]
to
\[
14 \times 5 = 70
\]
Therefore, the theoretical range of variable Y was 14--70.

\subsection{Calculation 1 --- Theoretical Score Range}

\[
X = 15\text{--}75 \qquad Y = 14\text{--}70
\]

This calculation helps us understand the minimum and maximum scores that
could theoretically be obtained by a respondent.

\begin{table}[htbp]
  \centering
  \singlespacing
  \caption{Descriptive Statistics}
  \label{tab:descriptive}
  \begin{tabularx}{\linewidth}{Xcccccc}
    \toprule
    Variable & N & Mean & Median & SD & Min & Max \\
    \midrule
    Gadget Use (X) & 69 & 56.45 & 58 & 9.09 & 31 & 74 \\
    Learning Motivation \& Social Interaction (Y) & 69 & 50.39 & 50 & 8.65 & 27 & 70 \\
    \bottomrule
  \end{tabularx}
\end{table}

Mean represents the average score of the respondents. Median represents
the middle value after all observations are arranged in order, while
standard deviation indicates how much the data are spread around the mean.

For example, the mean of X is 56.44927536, meaning that the average
gadget-use score among the 69 students was approximately 56.45. Similarly,
the mean of Y was 50.39130435, meaning that the average score for learning
motivation and social interaction was approximately 50.39.

This section is important because descriptive statistics provides an
initial overview of the data before we examine relationships between
variables.
